\properlane{word-heals}{The word heals and teaches faith}

The father's first request asks Christ to travel; Christ's answer asks the
father to travel. That reversal is the Gospel's decisive movement. The
petitioner wants the healer beside the child. Instead, he is given a word
on which to act while healer and child remain apart. His going home is
neither resignation nor abandonment of the sick. It is the conduct of
someone who has begun to accept a power greater than the means he first
imagined necessary. The word that heals also teaches its hearer how to
believe.

\subsection{A request corrected without being rejected}
Gregory begins with a question that the story itself provokes. Why does
Christ say that the petitioner does not believe unless he sees signs,
when the petitioner has come precisely because he believes Christ can
save his son? Gregory refuses to treat the request as devoid of faith.
The father would not seek healing from someone he thought powerless.
But he seeks bodily presence from the Lord who is nowhere absent in his
divine power. Gregory locates the defect within the belief: it acknowledges
a healer while limiting the reach of his healing.\footnote{Gregory,
\work{Homilies on the Gospels}, II.28.1.}

This reading makes the repeated request intelligible. The father is not
delivering a theological proposition about divine omnipresence. He is
trying to save a child before death comes, and his urgency clings to the
means he understands. Christ grants the desired life while removing the
restriction placed on his power. Gregory's Creator heals by command;
bodily distance does not isolate the sick child from him. The mercy and
the correction are one act. The father receives more than a successful
intervention: his estimate of the one who intervenes has to change.

Augustine approaches the same question through the preceding Samaritan
encounter. The Samaritans heard and believed; the Galilean reception is
marked by signs already seen. He takes Christ's rebuke as a disclosure
of the petitioner's heart, reading the initial approach as lukewarm,
cold, or not yet believing. The narrative's final declaration of belief
supports that severe judgment for him.\footnote{Augustine,
\work{Tractates on John}, XVI.1--5.}
Gregory's imperfect believer and Augustine's initially unbelieving
petitioner are different accounts. Yet the difference does not undo
their common Christological point: the Lord cannot be measured by the
petitioner's demand for a visible journey.

Augustine compares both the ruler and Thomas with those who believe
without demanding such sight. The comparison praises trust in Christ's
word; it does not remove the Gospel's testimony. The father still asks
the servants about the hour, and John records the answer. The confirmation
is received after a word has already become the ground of action. The
distinction is between evidence welcomed as testimony and a visible
condition imposed on the Lord before one will obey him. It is not a
choice between faith and careful attention to what happened.

\subsection{Confession frees the petitioner from self-sufficiency}
The Introit and Collect prepare an assembly capable of hearing this
Gospel. The entrance confession names disobedience; the Collect asks
for pardon and peace. Both address people who cannot repair their condition
simply by declaring themselves secure. A request for healing can begin
with an honest account of need. Its confidence rests on God's mercy,
not on the claim that the petitioner has already mastered himself.

Jerome's explanation of Azarias protects that honesty from becoming
cruelty. The faithful young men speak for their people; their confession
does not make their personal wrongdoing the cause of the furnace.
Likewise, the ruler's child has not been accused of causing his illness.
Corporate penitence and compassionate petition can stand together
without turning suffering into a public verdict on the sufferer.
\footnote{Jerome, \work{Commentary on Daniel}, I, at 3:29.}
The Introit gives the assembly a truthful first-person voice for repentance.
It gives no observer a warrant to explain someone else's pain by an
unrevealed offence.

The Collect's secure mind then differs from the demand to control every
event. Schuster describes pardon as restoring a heart capable of
peaceful service. That service is possible because guilt and inner
division no longer have the final word. In the Gospel the father's
departure supplies a concrete image of such released action: he can
move because Christ has spoken. He does not first receive possession of
the whole outcome.\footnote{Schuster, \work{The Sacramentary}, III,
pp.~175--177.}
The moral force of the scene is therefore precise. Trust permits the next
faithful action; it does not confer foreknowledge or require passivity.

Ephesians gives this action its ordinary texture. Walking wisely requires
discernment amid evil days, restraint rather than drunkenness, attention
rather than empty song, and reciprocal service rather than insistence
on precedence. Chrysostom's exposition makes prudence costly: one may
have to lose an advantage rather than return hostility with hostility.
Belief in the word is not confined to a dramatic moment of danger. It
becomes a pattern of choices in which the Lord's will matters more than
the immediate satisfaction of anger or appetite.

Gregory adds an especially sharp test. Why is Jesus willing to go to a
centurion's servant in the comparison Gospel, but does not go to this
ruler's son? Gregory hears a rebuke of the human habit of honoring wealth
and position rather than the person made in God's image.\footnote{Gregory,
II.28.2--3. The centurion's servant is Gregory's comparison, not the
patient in John~4.}
The healing word loosens not only a mistaken idea of distance but a
mistaken ordering of human worth. Faith in the Lord whose power reaches
everywhere cannot make a low-ranking neighbor beneath one's attention.

\subsection{What if healing does not arrive at once?}
The particular child is healed, and the Gospel gives the corresponding
hour. The rest of the Mass prevents the reader from converting that
particular mercy into a promise of immediate relief for every request.
The Gradual speaks of due season; the Offertory continues to weep; the
Communion seeks comfort in humiliation. An assembly can sing all these
words while containing people whose wounds remain.

Augustine's physician image in Psalm~144 directly addresses the person
whose good request seems unanswered. A physician's care includes judgment
about when and what to give. Desire and benefit need not coincide at the
moment the patient chooses. Augustine therefore urges the petitioner
not to lose heart. But he carries the argument beyond waiting for the
same earthly benefit. He asks the one grateful for health, family, and
other gifts to desire God himself. God's own gift of himself surpasses
the benefits through which the petitioner first learned his goodness.
\footnote{Augustine, Ps.~144, \S\S14--17.}

The progression matters. If the physician image ended with a guaranteed
later bodily recovery, it would still make divine goodness depend on
the eventual fulfilment of the petitioner's original plan. Augustine
instead directs hope toward salvation's final disclosure. This does not
make the son unimportant or turn his restored life into a disposable
illustration. His life is a real good, and his household rejoices in a
real mercy. The greater hope lets those who receive a different temporal
answer remain within the same prayer of trust.

The Offertory teaches a related freedom. Remembering Zion does not make
Babylon pleasant. Augustine's exiles sit by the currents and weep; they
do not need to deny their captivity to keep their desire directed toward
home. A sufferer can grieve without deciding that God's promise is void.
The object of hope holds beyond the circumstances that give rise to tears.
That is why the Communion can speak consolation in humiliation rather
than demand humiliation's prior disappearance.

\subsection{Heavenly medicine and the promised Word}
The Secret gives the language of healing a sacramental location. It asks
these mysteries to supply heavenly medicine and expel the vices of the
heart. The prayer does not retell the child's fever as though it were
merely an allegory for sin. It names another real need: the heart that
receives mercy also needs freedom from what deforms its love. Schuster's
exposition presents the Eucharist as an effective remedy for this
condition.\footnote{Schuster, III, p.~177.}
The bodily healing in John and the heart's healing in the Secret meet in
the saving power of Christ without becoming interchangeable ailments.

The Gradual's nourishment now acquires a sacramental horizon. Schuster
follows the food given in due season through the gifts that sustain faith
to the Eucharist and finally to the unveiled sight of God. Reception under
the veil of faith belongs to a journey, not to a failure of divine generosity.
The Alleluia's prepared heart answers that gift with praise. In Schuster's
reading God is the soul's glory: the healed life finds its measure in him,
not in the prestige that Gregory exposes as misleading.

At Communion, Augustine's account of remembrance resolves a potential
difficulty. Why ask an all-knowing God to remember? Because the servant
is asking for performance of the promise, not correction of the divine
memory. Hope itself has been given through the promise. The prayer can
therefore insist without accusing God of forgetfulness.
\footnote{Augustine, Ps.~118, \S\S50--51.}
Schuster's Eucharistic exposition identifies the Word with Christ, in
whom the Father's delight and the faithful's hope meet. His liturgical
reading intensifies the personal character of the promise; Augustine's
exegesis supplies the meaning of the petition that addresses it.

The Postcommunion completes the movement by asking for obedience. At
first this may sound like a reversal: why request worthiness after
receiving a remedy? The prayer's own grammar answers. The Lord is asked
to make the faithful obey. The disposition is itself sought as grace.
Augustine, continuing his exposition of Psalm~118, rejects the boast
of a righteousness treated as one's independent possession. Schuster
describes habitual docility as the disposition for fruitful reception.
Together they make possible a life that takes the commandments seriously
without treating them as a price paid to acquire God's generosity.
\footnote{Augustine, Ps.~118, \S56; Schuster, III, pp.~177--178.}

\subsection{The incarnate healer and enduring peace}
In the commemorative branch, the Maternity Collect identifies the healer
as the Word who took flesh from Mary. Augustine's concluding reflection
on John~4 holds together the Creator and his real human birth. The Son
is not less divine because he comes into the world he made. Nor does
his power at a distance make his human body unreal. The petition for
Mary's intercession belongs within this confession of the incarnate Word.

The Maternity Secret asks for both present and perpetual peace; its
Postcommunion asks for cleansing and participation in heaven's remedy.
These prayers widen the petitioner's horizon from the immediate need
toward enduring communion. They leave the source of healing with God.
The Trinity Preface, which remains the same with or without the
commemoration, prevents the healing story from ending in fascination
with power alone. The assembly thanks and adores the Father, Son, and
Holy Spirit in one divine majesty. Faith's end is communion with the
living God whom the gifts reveal.

\begin{fourSenses}
\item[Literal.]
A father seeks help for a dying child, accepts Jesus' word, and learns
that the cure coincided with it. The other passages retain their own
subjects: Azarias's communal confession, wise conduct, creaturely food,
exile, and a servant's hope. The prayers ask pardon, medicine, peace,
and obedience; the bodily cure remains bodily.
\item[Allegorical.]
Gregory's Creator heals by his command, and Augustine's incarnate Word
enters the people he made. Christ gathers the believing household and
nourishes his Church. Schuster's Eucharistic reading joins the promised
Word and heavenly medicine in the gift of Christ himself.
\item[Moral.]
Trust takes the next obedient step without claiming mastery of the
outcome. It repents, attends to the word, refuses rank-based contempt,
and serves others. The faithful ask for the grace of obedience; lament,
care for the sick, and inquiry into what happened remain compatible
with that trust.
\item[Anagogical.]
The physician's final gift exceeds temporary relief. Augustine's hope
in resurrection and Schuster's movement from sacramental nourishment
to vision direct desire toward God himself. Present and perpetual peace
belong to one hope whose fulfilment is not exhausted by any earthly cure.
\end{fourSenses}
